\section{ModelNet40: Classification and Geometric Fidelity}
\label{sec:modelnet-results}

OA and mAcc below are the recorded batch-aggregated statistics defined
in Equation~\eqref{eq:reported-classification}. All result values are
preserved from the HPC runs. This distinction matters when interpreting
small differences or comparing against externally reported accuracies.

\subsection{Densifying the Original Observation}
\label{sec:linea-results}

None of the five Line~A upsampled representations exceeded the original
1,024-point baseline. Table~\ref{tab:mn-a} reports the selected-checkpoint
overall accuracy (Best OA), epoch-200 accuracy (Final OA), and mean class
accuracy at the selected checkpoint. PU-GCN was closest in Best OA,
at 91.63\% compared with 91.95\%, while PDANS was nearly identical
at 91.62\%. The largest deficit was 1.41 percentage points (pp) for
PU-EdgeFormer. Each representation had its own classifier trained from random
initialisation, so these losses occurred even when the downstream model was
trained on the generated representation. Each row represents one
seed-42 classifier training run over 200 epochs.

\begin{table}[t]\centering
\setlength{\tabcolsep}{4pt}
\begin{tabular}{lrrrrr}\toprule
Representation & Points & Best OA & Final OA & mAcc & $\Delta$ OA \\
\midrule\tableinput{tables/modelnet_classification_A.tex}\bottomrule
\end{tabular}
\caption{ModelNet40 Line~A classification, with accuracies in percent.
$\Delta$ OA is the Best-OA change from Original~1024 in percentage points.}\label{tab:mn-a}
\end{table}

The same qualitative conclusion held at the fixed final epoch. The original
representation reached 91.31\%, compared with 91.00\% for PU-GCN and
90.87\% for PDANS. Thus the negative Line~A outcome was not produced only
by reporting the maximum of a noisy trajectory. The differences remain descriptive because
the test set was also used for checkpoint selection and training was not
repeated across independent random seeds.

Mean class accuracy and overall accuracy did not always move together.
PDANS and EAR had slightly higher reported mAcc than the original baseline
despite lower OA. Reported OA averages batch accuracies, whereas reported
mAcc averages the class-specific batch ratios across categories. Their
weights therefore differ, and a gain in one aggregate need not produce
a gain in the other. Without retained object predictions, a sample-weighted
recalculation or attribution to particular misclassified objects is not
possible. The category scores supplied with the runs are
analysed in Section~\ref{sec:mn-categories}. They are not converted into
integer correct-prediction counts because per-object predictions were not
retained.

\subsection{Recovery after Fourfold Sparsification}
\label{sec:lineb-results}

Only PU-Net exceeded the native 256-point Line~B baseline in Best OA:
91.27\% versus 90.85\%, a gain of 0.42~pp
(Table~\ref{tab:mn-b}). It remained 0.68~pp below Original~1024.
Restoring the original point count therefore recovered only part of the
classification gap. The other methods performed below the sparse baseline,
with deficits from 0.51~pp for PDANS to 2.04~pp for EAR.

\begin{table}[t]\centering
\setlength{\tabcolsep}{4pt}
\begin{tabular}{lrrrrr}\toprule
Representation & Points & Best OA & Final OA & mAcc & $\Delta$ OA \\
\midrule\tableinput{tables/modelnet_classification_B.tex}\bottomrule
\end{tabular}
\caption{ModelNet40 Line~B classification after $1,024\rightarrow256$
sparsification and $256\rightarrow1,024$ upsampling. The immediate baseline
retains its native 256 points. Definitions match Table~\ref{tab:mn-a}.}
\label{tab:mn-b}\end{table}

At epoch~200, PU-Net retained a smaller advantage of 0.19~pp over the
sparse baseline, 90.65\% versus 90.46\%. The recorded means over epochs
181--200 were 90.85\% and 90.32\%, respectively. This late-training
record supports a persistent descriptive advantage rather than an isolated
best checkpoint. Epochs from one optimisation trajectory are correlated,
however, and their variation cannot substitute for training-seed uncertainty.

Figure~\ref{fig:mn-training} shows the complete test-accuracy trajectories.
The later epochs show the sustained separation of EAR and PU-EdgeFormer
below the sparse baseline. The best-to-final gap reached 0.88~pp for
Line~B PDANS and 0.79~pp for EAR. Reporting only Best OA would obscure
that sensitivity. Under this training procedure, PU-Net provided the most
useful recovery representation among the tested methods; the evidence does
not establish a statistically significant or architecture-independent gain.

\thesisfigure{mn_training}{Classification during training}{Test OA at every
recorded epoch for the canonical branches. No smoothing or averaging across
runs is applied. The curves describe one seed per representation.}{fig:mn-training}

\thesisfigure{mn_late_training}{Later classification trajectories}{Unsmoothed
test OA from epochs 100--200. This complements the complete trajectories
by resolving the small late-training separations. Each curve is one
training run; no seed confidence interval is implied.}{fig:mn-late}
Figure~\ref{fig:mn-late} resolves the later training behaviour that is
compressed in a full-epoch plot. In Line~A, the final train--test gap
was 6.51~pp for Original and 8.23--9.38~pp for the upsampled inputs.
In Line~B it was 2.52~pp for Sparse and 7.24--9.58~pp after upsampling.
Higher training accuracy therefore coexisted with weaker test accuracy
for most generated representations. This observation is consistent with
greater fitting to the training representation, but does not isolate
overfitting from representation quality or optimisation effects.

\FloatBarrier
\subsection{Category-Dependent Classification Response}
\label{sec:mn-categories}

The overall results conceal different category responses. The exported
Line~A class scores contain both gains and losses for every method.
PU-GCN improved 18 categories, left 5 unchanged, and reduced 17; PDANS
improved 16 and reduced 18. Thus a method can lose overall accuracy
without degrading every category. These are descriptive comparisons of
the retained class aggregates, not paired counts of correctly classified
objects.

Table~\ref{tab:mn-category-counts} summarises all 40 categories in each
method--line pair. In Line~B, PU-Net had the most favourable balance:
15 improved, 15 unchanged, and 10 reduced. EAR reduced 24 categories,
while PU-EdgeFormer reduced 20. This broader distribution of losses
helps interpret their lower overall accuracy; it does not identify
which individual objects changed their predictions.

\begin{table}[!t]\centering
\begin{tabular}{llrrrr}\toprule
Line & Method & Improved & Same & Reduced & Mean $\Delta$ (pp)\\\midrule
A & EAR & 16 & 9 & 15 & +0.19 \\
A & PDANS & 16 & 6 & 18 & +0.44 \\
A & PU-Net & 15 & 6 & 19 & -0.54 \\
A & PU-GCN & 18 & 5 & 17 & -0.26 \\
A & PU-EdgeFormer & 17 & 7 & 16 & -1.66 \\
B & EAR & 8 & 8 & 24 & -2.86 \\
B & PDANS & 10 & 12 & 18 & -1.17 \\
B & PU-Net & 15 & 15 & 10 & +0.49 \\
B & PU-GCN & 13 & 8 & 19 & -0.96 \\
B & PU-EdgeFormer & 11 & 9 & 20 & -2.15 \\
\bottomrule\end{tabular}
\caption{Direction of reported category-score changes relative to the
native baseline. Each row covers all 40 categories. The mean is an
unweighted mean of archived category deltas and is not reconstructed OA.}
\label{tab:mn-category-counts}\end{table}

Figures~\ref{fig:mn-cat-A-1}, \ref{fig:mn-cat-A-2}, \ref{fig:mn-cat-B-1}, and~\ref{fig:mn-cat-B-2} retain the complete
category distribution using a common colour range. Categories are split
alphabetically across two panels per line, rather than selected by the
size or sign of a result. For example, Line~B PU-Net gained 25.00~pp on
door but lost 22.62~pp on curtain. Its positive aggregate result therefore
coexists with a substantial category-specific failure. Line~B PU-GCN
gained 11.11~pp on flower pot but lost 16.67~pp on bench. These differences motivate retaining the complete category distribution
in the heatmaps; the underlying scores remain in the archived export.

\thesisfigure{mn_categories_A_1}{Line A category changes, part 1}{Reported class-score changes for Line~A, alphabetical categories 1--20. Blue denotes improvement and red denotes reduction from the native baseline; the colour range is shared by all four maps. PU-EF denotes PU-EdgeFormer. These are archived category aggregates, with no reconstructed object-level counts.}{fig:mn-cat-A-1}
\thesisfigure{mn_categories_A_2}{Line A category changes, part 2}{Reported class-score changes for Line~A, alphabetical categories 21--40. Blue denotes improvement and red denotes reduction from the native baseline; the colour range is shared by all four maps. PU-EF denotes PU-EdgeFormer. These are archived category aggregates, with no reconstructed object-level counts.}{fig:mn-cat-A-2}
\thesisfigure{mn_categories_B_1}{Line B category changes, part 1}{Reported class-score changes for Line~B, alphabetical categories 1--20. Blue denotes improvement and red denotes reduction from the native baseline; the colour range is shared by all four maps. PU-EF denotes PU-EdgeFormer. These are archived category aggregates, with no reconstructed object-level counts.}{fig:mn-cat-B-1}
\thesisfigure{mn_categories_B_2}{Line B category changes, part 2}{Reported class-score changes for Line~B, alphabetical categories 21--40. Blue denotes improvement and red denotes reduction from the native baseline; the colour range is shared by all four maps. PU-EF denotes PU-EdgeFormer. These are archived category aggregates, with no reconstructed object-level counts.}{fig:mn-cat-B-2}

\FloatBarrier



\subsection{Genuine-Sampling Controls and Point Reception}
\label{sec:controls}

Increasing genuine surface samples also had a limited return above 1,024
points. Mesh-ref~256 reached 90.88\% Best OA, close to the decimated
baseline's 90.85\%. Mesh-ref~4096 reached 92.00\%, only 0.05~pp
above Original~1024. At epoch~200 it was below Original~1024,
91.18\% versus 91.31\% (Figure~\ref{fig:mn-controls}). These controls
suggest diminishing benefit from additional surface samples for the configured
PointNet++ SSG classifier, rather than a general upper bound on density.

Four training shapes in the mesh-reference builds were reused instead of
regenerated under the final stable-seed rule. No test object was affected,
so the equal-cardinality test geometry remains usable. The two trained
mesh-reference controls are nevertheless provisional: their small OA
differences should not support a strong causal claim without rebuilding
those training shapes and repeating the control training.

\thesisfigure{mn_controls}{Genuine density controls}{Selected and final OA for
genuine mesh samples and the native sparse baseline. Four reused training
shapes limit the mesh-reference training comparison; the test geometry
contains no such exception.}{fig:mn-controls}

The fixed-radius grouping and finite neighbourhoods of PointNet++ \cite{qi2017pointnetpp} motivate a possible explanation for the limited return. In the retained reception probe, the genuine and
decimated 256-point inputs required repeated indices in 59.2\% of
neighbourhood slots. Original~1024 discarded 31.9\% of points found
within its queried neighbourhoods, while Mesh-ref~4096 discarded 78.7\%.
These are local grouping statistics with a fixed neighbourhood budget,
not percentages of globally unique points removed from the cloud.

Figure~\ref{fig:mn-sa1} reports repetition and discard rates. A fourfold
increase in file cardinality did not produce a fourfold increase in distinct
local evidence passed to the next layer. This is consistent with saturation
under fixed grouping radii and capacities. Reception is not a uniquely
identified cause of the accuracy differences: shape deformation, feature
extraction, and optimisation can also contribute.

\thesisfigure{mn_sa1}{PointNet++ neighbourhood reception}{Recorded first-layer
discard and repeated-slot rates for the retained representations. Rates are
local grouping statistics, not complementary fractions of an entire cloud.
The architecture is specified in Section~\ref{sec:classifier-setup}.}{fig:mn-sa1}

The earlier 512-point controls further show that the input protocol matters.
PDANS reached 91.47\% Best OA from 512 visible points against a native
baseline of 91.26\%, whereas the final 256-point comparison lost
0.51~pp. Both visible support and downstream input size changed;
this is not an isolated estimate of sparsity. The complete earlier
runs remain in the archived HPC results rather than being pooled
with the final two-line comparison.

\subsection{Equal-Cardinality Geometry}
\label{sec:equaln-geometry}
\label{sec:geometry-audit}

PU-GCN achieved the lowest mean unsquared Chamfer distance in both lines.
The comparison used 2,468 paired test objects, a 4,096-point independent
mesh reference in Line~A, and Original~1024 in Line~B.
Tables~\ref{tab:geom-a} and~\ref{tab:geom-b} give the means and archived
object-bootstrap intervals from 4,000 resamples of the 2,468 objects,
using the resampling principle of Efron \cite{efron1979bootstrap}. These intervals describe variation across test
objects, not training seeds or checkpoints.

\begin{table}[t]\centering\setlength{\tabcolsep}{4pt}
\begin{tabular}{lrrrrr}\toprule
Method & CD & CD 95\% interval & HD & $|\Delta\mathrm{NUC}|$ & P2F$^{*}$ \\
\midrule\tableinput{tables/modelnet_geometry_A.tex}\bottomrule
\end{tabular}
\caption{Line~A geometry against the independent 4,096-point reference.
CD and HD are unsquared distances in normalised coordinates.
$^{*}$P2F uses independently normalised clouds and meshes.}\label{tab:geom-a}
\end{table}

\begin{table}[t]\centering\setlength{\tabcolsep}{4pt}
\begin{tabular}{lrrrrr}\toprule
Method & CD & CD 95\% interval & HD & $|\Delta\mathrm{NUC}|$ & P2F$^{*}$ \\
\midrule\tableinput{tables/modelnet_geometry_B.tex}\bottomrule
\end{tabular}
\caption{Line~B geometry against Original~1024. Definitions match
Table~\ref{tab:geom-a}. The sparse 256-point baseline has a separate
256-point mesh reference and is not ranked in this table.
The P2F column retains the separate cloud-to-mesh convention.}\label{tab:geom-b}
\end{table}

The Chamfer advantage did not imply the best result under every metric.
PU-GCN also had the lowest mean Line~A HD, but PU-Net had the lowest
Line~B HD: 0.1083 compared with 0.1514 for PU-GCN.
Figure~\ref{fig:mn-geometry} shows the change in ordering. Chamfer distance
averages nearest-neighbour residuals, whereas HD records the largest gap
or outlier for each object. The rankings indicate a trade-off between average residuals and extreme
local residuals under the retained normalised-reference protocol. Equal
point counts control cardinality, but the independently normalised mesh
reference can still contribute translation and scale differences
(Section~\ref{sec:fc-4-10-1}). These results are not presented as a
common-transform surface-error ranking.

\thesisfigure{mn_geometry}{Geometric fidelity at equal cardinality}{Mean CD
and HD across 2,468 test objects. Error bars are archived 95\%
object-bootstrap intervals. The reference differs between the two lines
as specified in Tables~\ref{tab:geom-a} and~\ref{tab:geom-b}.}{fig:mn-geometry}

Uniformity also requires a reference. A low NUC indicates less variation in
local counts, but an object surface need not have the lowest achievable
NUC. The analysis therefore averages each object's absolute deviation from
its reference NUC; this differs from the absolute difference between dataset
means. Query centres were not fully shared by shape ID across methods,
so NUC differences include query-placement variation and remain secondary
to the principal distance and task results.

The P2F columns in Tables~\ref{tab:geom-a} and~\ref{tab:geom-b}
retain the completed secondary surface audit. Clouds and meshes were
independently centred and scaled, so even the original mesh-sampled
input had mean P2F 0.08662; the sparse control had 0.08675.
PU-Net had the lowest Line~A mean P2F (0.08413), and PDANS the lowest
Line~B value (0.08365). This ordering differs from CD and HD, but the
coordinate convention prevents interpreting it as an absolute
common-frame surface ranking. Co-location in the tables does not
make the reference protocols interchangeable.

The mean geometric ranking is supported by the number of objects on which
a method wins, rather than by a few extreme values alone.
PU-GCN had the lowest CD on 2,319 of 2,468 Line~A objects (93.96\%)
and 2,303 Line~B objects (93.31\%). For HD, however, its winner share
changed from 72.65\% in Line~A to 9.60\% in Line~B.
PU-Net had the lowest HD on 2,218 Line~B objects (89.87\%).
Figure~\ref{fig:mn-winners} therefore exposes a systematic disagreement
between average nearest-neighbour fidelity and extreme-error behaviour.

\thesisfigure{mn_object_winners}{Object-wise geometry winners}{Share of
2,468 test objects on which each method has the lowest archived CD or HD.
All five methods and both lines are retained. Winner counts are taken
from the HPC object-wise summary; they measure geometric comparisons,
not classification successes.}{fig:mn-winners}



\FloatBarrier
\subsection{Multi-Method Point-Cloud Evidence}
\label{sec:modelnet-visual}

Figures~\ref{fig:mn_atlas_A_airplane} and~\ref{fig:mn_atlas_B_chair}
retain two of the eight archived HPC examples, one per comparison line.
The airplane illustrates agreement between the lowest CD and HD;
the chair illustrates their different method orderings. These examples
were selected for this geometric contrast, not to estimate its frequency.
No per-object classification transition is assigned because complete
classifier predictions were not retained.

Each plate contains seven panels: input, reference and all five methods.
Every output point in panels (c--g) is coloured by its Euclidean
nearest-neighbour distance to the same reference, with one full-range
colour scale shared across methods within that plate. All stored points
are shown under the same orthographic view, extent and marker size.
This colour measures output-to-reference deviation only; CD also
includes reference-to-output coverage, and HD takes the worst of both
directions. Colour intensity is consequently not the complete CD or HD.

For airplane\_0627 in Line~A, PU-GCN had the lowest CD (0.02979) and
HD (0.09912), while PU-Net had 0.04231 and 0.11862, respectively.
PU-EdgeFormer's HD was close to PU-GCN's at 0.09976 despite its higher
CD of 0.03285. Figure~\ref{fig:mn_atlas_A_airplane} locates deviations
around the same wings, tail and fuselage under the retained reference.

\thesisfigure{mn_atlas_A_airplane}{Airplane multi-method comparison, Line A}{
HPC airplane\_0627. (a) Input, 1,024 points; (b) independent mesh
reference, 4,096 points; (c--g) five 4,096-point outputs, all coloured
by output-to-reference Euclidean distance in normalised coordinates.
PU-EF denotes PU-EdgeFormer. Every method shares the view and colour
range; the range is specific to this object.}{fig:mn_atlas_A_airplane}

For chair\_0890 in Line~B, PU-GCN had lower CD than PU-Net, 0.03281
versus 0.04330, but higher HD, 0.18622 versus 0.06956.
Figure~\ref{fig:mn_atlas_B_chair} therefore complements the dataset-level
winner shares: better average proximity can coexist with a worse
extreme residual for the same object. The sparse input, reference and
generated outputs retain the seat, back and narrow supports in a
common view. Neither their denser appearance nor their geometric
ordering establishes a classification improvement.

\thesisfigure{mn_atlas_B_chair}{Chair multi-method comparison, Line B}{
HPC chair\_0890. (a) Sparse input, 256 points; (b) Original~1024
reference; (c--g) five 1,024-point outputs, all coloured by
output-to-reference Euclidean distance in normalised coordinates.
PU-EF denotes PU-EdgeFormer. Every method shares the view and colour
range; the range is specific to this object.}{fig:mn_atlas_B_chair}
\FloatBarrier

